% ============================================================
% CHAPITRE 5 : IMPLÉMENTATION
% ============================================================

\chapter{Implémentation}

\section{Introduction}

Ce chapitre détaille l'implémentation logicielle du système DMS en Python, incluant l'entraînement des modèles CNN, les algorithmes de vision par ordinateur et la logique de traitement. Chaque composant est présenté avec son code source, les algorithmes associés et les équations utilisées.

\section{Environnement de Développement}

\subsection{Outils et Technologies}

\begin{table}[H]
\centering
\caption{Environnement de développement du système DMS.}
\label{tab:dev_env}
\renewcommand{\arraystretch}{1.3}
\begin{tabularx}{\textwidth}{|l|l|X|}
\hline
\rowcolor{stellantisblue!15}
\textbf{Outil} & \textbf{Version} & \textbf{Rôle} \\
\hline
Python & 3.10+ & Langage principal de développement \\
\hline
PyTorch & $\geq 2.0.0$ & Framework de deep learning pour les CNN \\
\hline
OpenCV & $\geq 4.8.0$ & Vision par ordinateur (détection, prétraitement) \\
\hline
NumPy & $\geq 1.24.0$ & Calcul numérique (matrices, algèbre linéaire) \\
\hline
Pandas & $\geq 2.0.0$ & Analyse et manipulation de données \\
\hline
Matplotlib & $\geq 3.7.0$ & Visualisation des résultats et graphiques \\
\hline
dlib & $\geq 19.24.0$ & Détection de landmarks faciaux \\
\hline
SciPy & $\geq 1.10.0$ & Calculs scientifiques (rotation, distances) \\
\hline
\end{tabularx}
\end{table}

\subsection{Structure du Projet}

\begin{lstlisting}[style=pythonstyle, caption={Structure du projet DMS.}, label={lst:project_structure}, language={}]
dms-project/
|-- src/
|   |-- __init__.py                 # Module principal
|   |-- video_acquisition.py        # Acquisition video
|   |-- face_eye_detection.py       # Detection visage/yeux
|   |-- head_pose.py                # Estimation pose tete
|   |-- behavioral_analysis.py      # Analyse comportementale
|   |-- ecu_decision.py             # Logique ECU
|   |-- scenarios.py                # Scenarios de test
|   |-- visualization.py            # Visualisation
|-- tests/
|   |-- test_dms.py                 # Tests unitaires (34 tests)
|-- matlab/
|   |-- dms_simulation.m            # Simulation MATLAB
|   |-- dms_simulink_model.m        # Modele Simulink
|   |-- simulink_model/
|       |-- ecu_stateflow.m         # Machine a etats
|-- main.py                         # Point d'entree
|-- requirements.txt                # Dependances
\end{lstlisting}

\section{Implémentation du Module d'Acquisition Vidéo}

\subsection{Classe VideoAcquisition}

Le module d'acquisition vidéo gère la capture et le prétraitement des images :

\begin{lstlisting}[style=pythonstyle, caption={Module d'acquisition vidéo.}, label={lst:video_acquisition}]
import cv2
import numpy as np

class VideoAcquisition:
    """Module d'acquisition video pour le systeme DMS."""
    
    def __init__(self, source=0, resolution=(640, 480), fps=30):
        """
        Initialise l'acquisition video.
        
        Args:
            source: Index camera (int) ou chemin fichier (str)
            resolution: Tuple (largeur, hauteur)
            fps: Images par seconde
        """
        self.source = source
        self.resolution = resolution
        self.fps = fps
        self.cap = None
    
    def open(self):
        """Ouvre la source video."""
        self.cap = cv2.VideoCapture(self.source)
        self.cap.set(cv2.CAP_PROP_FRAME_WIDTH, self.resolution[0])
        self.cap.set(cv2.CAP_PROP_FRAME_HEIGHT, self.resolution[1])
        self.cap.set(cv2.CAP_PROP_FPS, self.fps)
        return self.cap.isOpened()
    
    def read_frame(self):
        """Lit une image de la source video."""
        if self.cap is None:
            return None
        ret, frame = self.cap.read()
        return frame if ret else None
    
    def preprocess(self, frame):
        """
        Preprocesse l'image pour la detection.
        
        Pipeline:
        1. Conversion en niveaux de gris
        2. Egalisation d'histogramme (CLAHE)
        """
        gray = cv2.cvtColor(frame, cv2.COLOR_BGR2GRAY)
        clahe = cv2.createCLAHE(
            clipLimit=2.0, 
            tileGridSize=(8, 8)
        )
        enhanced = clahe.apply(gray)
        return enhanced
    
    def generate_frames(self):
        """Generateur d'images preprocessees."""
        while True:
            frame = self.read_frame()
            if frame is None:
                break
            gray = self.preprocess(frame)
            yield frame, gray
    
    def release(self):
        """Libere les ressources."""
        if self.cap is not None:
            self.cap.release()
\end{lstlisting}

\subsection{Algorithme CLAHE}

L'égalisation adaptative d'histogramme à contraste limité (\textit{CLAHE -- Contrast Limited Adaptive Histogram Equalization}) divise l'image en tuiles et applique une égalisation locale :

\begin{equation}
H_{\text{CLAHE}}(v, r) = \text{round}\left(\frac{\text{CDF}_r(v) - \text{CDF}_{r,\min}}{(M_r \times N_r) - \text{CDF}_{r,\min}} \times (L - 1)\right)
\label{eq:clahe}
\end{equation}

où $r$ désigne la tuile locale et les CDF sont écrêtées au \textit{clipLimit} pour éviter l'amplification du bruit.

\section{Implémentation des CNN}

\subsection{EyeStateCNN -- Architecture Complète}

\begin{lstlisting}[style=pythonstyle, caption={Architecture complète du EyeStateCNN.}, label={lst:eye_cnn}]
import torch
import torch.nn as nn
import torch.nn.functional as F

class EyeStateCNN(nn.Module):
    """
    Reseau de neurones convolutif pour la classification
    de l'etat des yeux (ouvert/ferme).
    
    Entree: Image 1x24x24 (niveaux de gris)
    Sortie: 2 classes (ouvert, ferme)
    """
    
    def __init__(self):
        super().__init__()
        # Bloc convolutif 1: 1 -> 32 filtres
        self.conv1 = nn.Conv2d(1, 32, kernel_size=3, padding=1)
        self.bn1 = nn.BatchNorm2d(32)
        
        # Bloc convolutif 2: 32 -> 64 filtres
        self.conv2 = nn.Conv2d(32, 64, kernel_size=3, padding=1)
        self.bn2 = nn.BatchNorm2d(64)
        
        # Bloc convolutif 3: 64 -> 128 filtres
        self.conv3 = nn.Conv2d(64, 128, kernel_size=3, padding=1)
        self.bn3 = nn.BatchNorm2d(128)
        
        # Couches communes
        self.pool = nn.MaxPool2d(2, 2)
        self.dropout = nn.Dropout(0.5)
        
        # Couches entierement connectees
        # Apres 3 couches de conv+pool: 24->12->6->3
        # Taille: 128 * 3 * 3 = 1152
        self.fc1 = nn.Linear(128 * 3 * 3, 256)
        self.fc2 = nn.Linear(256, 2)
    
    def forward(self, x):
        """
        Propagation avant.
        
        Args:
            x: Tensor de forme (batch, 1, 24, 24)
        Returns:
            Tensor de forme (batch, 2) - logits
        """
        # Bloc 1: Conv -> BN -> ReLU -> MaxPool
        # (B, 1, 24, 24) -> (B, 32, 12, 12)
        x = self.pool(F.relu(self.bn1(self.conv1(x))))
        
        # Bloc 2: Conv -> BN -> ReLU -> MaxPool
        # (B, 32, 12, 12) -> (B, 64, 6, 6)
        x = self.pool(F.relu(self.bn2(self.conv2(x))))
        
        # Bloc 3: Conv -> BN -> ReLU -> MaxPool
        # (B, 64, 6, 6) -> (B, 128, 3, 3)
        x = self.pool(F.relu(self.bn3(self.conv3(x))))
        
        # Aplatissement: (B, 128, 3, 3) -> (B, 1152)
        x = x.view(x.size(0), -1)
        
        # FC1 + ReLU + Dropout: (B, 1152) -> (B, 256)
        x = self.dropout(F.relu(self.fc1(x)))
        
        # FC2 (sortie): (B, 256) -> (B, 2)
        x = self.fc2(x)
        
        return x
\end{lstlisting}

\subsection{FaceDetectionCNN -- Architecture Complète}

\begin{lstlisting}[style=pythonstyle, caption={Architecture du FaceDetectionCNN.}, label={lst:face_cnn}]
class FaceDetectionCNN(nn.Module):
    """
    CNN pour la validation des detections de visage.
    
    Entree: Image 1x64x64 (niveaux de gris)
    Sortie: 2 classes (visage, non-visage)
    """
    
    def __init__(self):
        super().__init__()
        # 4 blocs convolutifs avec profondeur croissante
        self.conv1 = nn.Conv2d(1, 16, kernel_size=3, padding=1)
        self.bn1 = nn.BatchNorm2d(16)
        self.conv2 = nn.Conv2d(16, 32, kernel_size=3, padding=1)
        self.bn2 = nn.BatchNorm2d(32)
        self.conv3 = nn.Conv2d(32, 64, kernel_size=3, padding=1)
        self.bn3 = nn.BatchNorm2d(64)
        self.conv4 = nn.Conv2d(64, 128, kernel_size=3, padding=1)
        self.bn4 = nn.BatchNorm2d(128)
        
        self.pool = nn.MaxPool2d(2, 2)
        self.adaptive_pool = nn.AdaptiveAvgPool2d(1)
        self.dropout = nn.Dropout(0.5)
        
        self.fc1 = nn.Linear(128, 64)
        self.fc2 = nn.Linear(64, 2)
    
    def forward(self, x):
        """
        Propagation avant.
        
        Flux de donnees:
        (B,1,64,64) -> Conv1 -> (B,16,32,32) -> Conv2 -> 
        (B,32,16,16) -> Conv3 -> (B,64,8,8) -> Conv4 -> 
        (B,128,4,4) -> AdaptiveAvgPool -> (B,128) -> 
        FC1 -> (B,64) -> FC2 -> (B,2)
        """
        x = self.pool(F.relu(self.bn1(self.conv1(x))))
        x = self.pool(F.relu(self.bn2(self.conv2(x))))
        x = self.pool(F.relu(self.bn3(self.conv3(x))))
        x = self.pool(F.relu(self.bn4(self.conv4(x))))
        x = self.adaptive_pool(x)
        x = x.view(x.size(0), -1)
        x = self.dropout(F.relu(self.fc1(x)))
        x = self.fc2(x)
        return x
\end{lstlisting}

\subsection{Processus d'Entraînement du CNN}

\subsubsection{Fonction de Perte}

L'entraînement utilise la fonction de perte \keyword{Cross-Entropy} pour la classification binaire :

\begin{equation}
\mathcal{L}_{\text{CE}} = -\frac{1}{N} \sum_{i=1}^{N} \left[ y_i \log(\hat{y}_i) + (1 - y_i) \log(1 - \hat{y}_i) \right]
\label{eq:cross_entropy}
\end{equation}

où $y_i \in \{0, 1\}$ est l'étiquette vraie et $\hat{y}_i = \text{softmax}(z_i)_1$ la probabilité prédite pour la classe positive.

\subsubsection{Optimiseur Adam}

L'optimiseur \keyword{Adam} (\textit{Adaptive Moment Estimation}) est utilisé avec les règles de mise à jour suivantes :

\begin{align}
m_t &= \beta_1 m_{t-1} + (1 - \beta_1) g_t \label{eq:adam_m} \\
v_t &= \beta_2 v_{t-1} + (1 - \beta_2) g_t^2 \label{eq:adam_v} \\
\hat{m}_t &= \frac{m_t}{1 - \beta_1^t}, \quad \hat{v}_t = \frac{v_t}{1 - \beta_2^t} \label{eq:adam_bias} \\
\theta_t &= \theta_{t-1} - \frac{\eta}{\sqrt{\hat{v}_t} + \epsilon} \hat{m}_t \label{eq:adam_update}
\end{align}

avec les hyperparamètres : $\eta = 0.001$, $\beta_1 = 0.9$, $\beta_2 = 0.999$, $\epsilon = 10^{-8}$.

\subsubsection{Code d'Entraînement}

\begin{lstlisting}[style=pythonstyle, caption={Boucle d'entraînement du CNN.}, label={lst:training}]
def train_model(model, train_loader, val_loader, 
                epochs=50, lr=0.001):
    """
    Entraine le modele CNN.
    
    Args:
        model: Modele PyTorch
        train_loader: DataLoader d'entrainement
        val_loader: DataLoader de validation
        epochs: Nombre d'epoques
        lr: Taux d'apprentissage
    """
    criterion = nn.CrossEntropyLoss()
    optimizer = torch.optim.Adam(model.parameters(), lr=lr)
    scheduler = torch.optim.lr_scheduler.ReduceLROnPlateau(
        optimizer, mode='min', patience=5, factor=0.5
    )
    
    best_val_acc = 0.0
    
    for epoch in range(epochs):
        # Phase d'entrainement
        model.train()
        train_loss = 0.0
        correct = 0
        total = 0
        
        for images, labels in train_loader:
            optimizer.zero_grad()
            
            # Propagation avant
            outputs = model(images)
            loss = criterion(outputs, labels)
            
            # Retropropagation
            loss.backward()
            optimizer.step()
            
            train_loss += loss.item()
            _, predicted = torch.max(outputs.data, 1)
            total += labels.size(0)
            correct += (predicted == labels).sum().item()
        
        train_acc = 100 * correct / total
        
        # Phase de validation
        model.eval()
        val_loss = 0.0
        val_correct = 0
        val_total = 0
        
        with torch.no_grad():
            for images, labels in val_loader:
                outputs = model(images)
                loss = criterion(outputs, labels)
                val_loss += loss.item()
                _, predicted = torch.max(outputs.data, 1)
                val_total += labels.size(0)
                val_correct += (predicted == labels).sum().item()
        
        val_acc = 100 * val_correct / val_total
        scheduler.step(val_loss)
        
        # Sauvegarde du meilleur modele
        if val_acc > best_val_acc:
            best_val_acc = val_acc
            torch.save(model.state_dict(), 'best_model.pth')
        
        print(f'Epoch [{epoch+1}/{epochs}] '
              f'Train Loss: {train_loss:.4f} '
              f'Train Acc: {train_acc:.2f}% '
              f'Val Acc: {val_acc:.2f}%')
\end{lstlisting}

\subsubsection{Rétropropagation du Gradient}

La rétropropagation calcule les gradients de la fonction de perte par rapport à chaque paramètre en utilisant la règle de la chaîne :

\begin{equation}
\frac{\partial \mathcal{L}}{\partial \mathbf{W}^{(l)}} = \frac{\partial \mathcal{L}}{\partial \mathbf{z}^{(l)}} \cdot \frac{\partial \mathbf{z}^{(l)}}{\partial \mathbf{W}^{(l)}} = \boldsymbol{\delta}^{(l)} \cdot (\mathbf{a}^{(l-1)})^T
\label{eq:backprop}
\end{equation}

où $\boldsymbol{\delta}^{(l)}$ est l'erreur rétropropagée à la couche $l$ et $\mathbf{a}^{(l-1)}$ les activations de la couche précédente.

Pour la couche de convolution :

\begin{equation}
\frac{\partial \mathcal{L}}{\partial \mathbf{W}_k^{(l)}} = \sum_{i} \mathbf{a}_i^{(l-1)} * \boldsymbol{\delta}_k^{(l)}
\label{eq:conv_backprop}
\end{equation}

\section{Implémentation de la Détection Visage/Yeux}

\subsection{Algorithme de Détection Complet}

\begin{algorithm}[H]
\SetAlgoLined
\caption{Détection du visage et des yeux}
\label{algo:detection}
\KwIn{Image BGR $\mathbf{I}$ du flux vidéo}
\KwOut{Liste de détections $\mathcal{D} = \{(bbox, \text{état\_yeux}, \text{confiance})\}$}

$\mathbf{I}_{\text{gray}} \leftarrow \text{cvtColor}(\mathbf{I}, \text{BGR2GRAY})$\;
$\mathbf{I}_{\text{eq}} \leftarrow \text{equalizeHist}(\mathbf{I}_{\text{gray}})$\;
$\mathcal{F} \leftarrow \text{HaarCascade}(\mathbf{I}_{\text{eq}}, \text{scaleFactor}=1.3, \text{minNeighbors}=5)$\;
$\mathcal{D} \leftarrow \emptyset$\;

\ForEach{$(x, y, w, h) \in \mathcal{F}$}{
    $\mathbf{R}_{\text{face}} \leftarrow \mathbf{I}_{\text{gray}}[y:y+h, x:x+w]$\;
    $\mathbf{R}_{\text{resize}} \leftarrow \text{resize}(\mathbf{R}_{\text{face}}, 64 \times 64)$\;
    $\mathbf{R}_{\text{norm}} \leftarrow \mathbf{R}_{\text{resize}} / 255.0$\;
    $c_{\text{face}} \leftarrow \text{softmax}(\text{FaceDetectionCNN}(\mathbf{R}_{\text{norm}}))[1]$\;
    
    \If{$c_{\text{face}} > \tau_{\text{face}}$}{
        $\mathcal{E} \leftarrow \text{HaarCascadeEyes}(\mathbf{R}_{\text{face}})$\;
        
        \ForEach{$(x_e, y_e, w_e, h_e) \in \mathcal{E}$}{
            $\mathbf{R}_{\text{eye}} \leftarrow \mathbf{R}_{\text{face}}[y_e:y_e+h_e, x_e:x_e+w_e]$\;
            $\mathbf{R}_{\text{eye,resize}} \leftarrow \text{resize}(\mathbf{R}_{\text{eye}}, 24 \times 24)$\;
            $\mathbf{R}_{\text{eye,norm}} \leftarrow \mathbf{R}_{\text{eye,resize}} / 255.0$\;
            $\text{état} \leftarrow \arg\max(\text{EyeStateCNN}(\mathbf{R}_{\text{eye,norm}}))$\;
            $c_{\text{eye}} \leftarrow \max(\text{softmax}(\text{EyeStateCNN}(\mathbf{R}_{\text{eye,norm}})))$\;
        }
        
        $\text{EAR} \leftarrow \frac{\|p_2 - p_6\| + \|p_3 - p_5\|}{2 \cdot \|p_1 - p_4\|}$\;
        $\mathcal{D} \leftarrow \mathcal{D} \cup \{(bbox, \text{état}, c_{\text{eye}}, \text{EAR})\}$\;
    }
}

\Return{$\mathcal{D}$}\;
\end{algorithm}

\subsection{Code d'Implémentation}

\begin{lstlisting}[style=pythonstyle, caption={Module de détection visage et yeux.}, label={lst:face_detection}]
class FaceEyeDetector:
    """Detecteur de visage et d'yeux combine."""
    
    def __init__(self):
        # Cascades de Haar (OpenCV)
        self.face_cascade = cv2.CascadeClassifier(
            cv2.data.haarcascades + 
            'haarcascade_frontalface_default.xml'
        )
        self.eye_cascade = cv2.CascadeClassifier(
            cv2.data.haarcascades + 
            'haarcascade_eye.xml'
        )
        
        # Modeles CNN
        self.face_cnn = FaceDetectionCNN()
        self.eye_cnn = EyeStateCNN()
        self.face_cnn.eval()
        self.eye_cnn.eval()
    
    def detect(self, frame, gray):
        """
        Detecte les visages et classifie l'etat des yeux.
        
        Args:
            frame: Image BGR originale
            gray: Image en niveaux de gris preprocessee
            
        Returns:
            Liste de detections avec confiance et etat
        """
        detections = []
        
        # Detection de visages par Haar
        faces = self.face_cascade.detectMultiScale(
            gray, scaleFactor=1.3, minNeighbors=5,
            minSize=(30, 30)
        )
        
        for (x, y, w, h) in faces:
            # Extraction ROI visage
            face_roi = gray[y:y+h, x:x+w]
            
            # Validation CNN du visage
            face_conf = self._classify_face(face_roi)
            
            if face_conf > 0.5:  # Seuil de confiance
                # Detection des yeux dans le visage
                eye_state, eye_conf = self._detect_eyes(
                    face_roi
                )
                
                detections.append({
                    'bbox': (x, y, w, h),
                    'face_confidence': face_conf,
                    'eye_state': eye_state,
                    'eye_confidence': eye_conf
                })
        
        return detections
    
    def _classify_face(self, face_roi):
        """Classifie la ROI comme visage/non-visage."""
        resized = cv2.resize(face_roi, (64, 64))
        tensor = torch.FloatTensor(resized).unsqueeze(0)
        tensor = tensor.unsqueeze(0) / 255.0
        
        with torch.no_grad():
            output = self.face_cnn(tensor)
            prob = F.softmax(output, dim=1)
        
        return prob[0][1].item()
    
    def _detect_eyes(self, face_roi):
        """Detecte et classifie l'etat des yeux."""
        eyes = self.eye_cascade.detectMultiScale(
            face_roi, scaleFactor=1.1, minNeighbors=3
        )
        
        if len(eyes) == 0:
            return 'unknown', 0.0
        
        states = []
        confidences = []
        
        for (ex, ey, ew, eh) in eyes[:2]:
            eye_roi = face_roi[ey:ey+eh, ex:ex+ew]
            resized = cv2.resize(eye_roi, (24, 24))
            tensor = torch.FloatTensor(resized).unsqueeze(0)
            tensor = tensor.unsqueeze(0) / 255.0
            
            with torch.no_grad():
                output = self.eye_cnn(tensor)
                prob = F.softmax(output, dim=1)
            
            state = 'open' if prob[0][0] > prob[0][1] \
                    else 'closed'
            conf = max(prob[0]).item()
            states.append(state)
            confidences.append(conf)
        
        # Vote majoritaire
        final_state = max(set(states), key=states.count)
        final_conf = np.mean(confidences)
        
        return final_state, final_conf
\end{lstlisting}

\section{Implémentation de l'Estimation de la Pose}

\begin{lstlisting}[style=pythonstyle, caption={Module d'estimation de la pose de la tête.}, label={lst:head_pose}]
class HeadPoseEstimator:
    """Estimateur de la pose de la tete par PnP."""
    
    # Points 3D du modele de visage generique (mm)
    MODEL_POINTS_3D = np.array([
        (0.0, 0.0, 0.0),         # Bout du nez
        (0.0, -330.0, -65.0),     # Menton
        (-225.0, 170.0, -135.0),  # Coin oeil gauche
        (225.0, 170.0, -135.0),   # Coin oeil droit
        (-150.0, -150.0, -125.0), # Coin bouche gauche
        (150.0, -150.0, -125.0),  # Coin bouche droit
    ], dtype=np.float64)
    
    def __init__(self, frame_width=640, frame_height=480):
        """
        Initialise l'estimateur.
        
        Calcul de la matrice intrinseque K approximee:
        fx = fy = largeur_image (approximation focale)
        cx, cy = centre de l'image
        """
        focal_length = frame_width
        center = (frame_width / 2, frame_height / 2)
        
        self.camera_matrix = np.array([
            [focal_length, 0, center[0]],
            [0, focal_length, center[1]],
            [0, 0, 1]
        ], dtype=np.float64)
        
        # Pas de distorsion (approximation)
        self.dist_coeffs = np.zeros((4, 1))
    
    def estimate(self, landmarks_2d):
        """
        Estime la pose de la tete.
        
        Args:
            landmarks_2d: Array (6, 2) des points 2D
            
        Returns:
            dict: {'yaw': float, 'pitch': float, 
                   'roll': float} en degres
        """
        landmarks = np.array(landmarks_2d, dtype=np.float64)
        
        # Resolution du probleme PnP
        success, rvec, tvec = cv2.solvePnP(
            self.MODEL_POINTS_3D,
            landmarks,
            self.camera_matrix,
            self.dist_coeffs,
            flags=cv2.SOLVEPNP_ITERATIVE
        )
        
        if not success:
            return {'yaw': 0.0, 'pitch': 0.0, 'roll': 0.0}
        
        # Conversion vecteur de Rodrigues -> matrice
        rotation_matrix, _ = cv2.Rodrigues(rvec)
        
        # Extraction des angles d'Euler
        angles = self._rotation_to_euler(rotation_matrix)
        
        return {
            'yaw': np.degrees(angles[1]),    # Lacet
            'pitch': np.degrees(angles[0]),  # Tangage
            'roll': np.degrees(angles[2])    # Roulis
        }
    
    def _rotation_to_euler(self, R):
        """
        Convertit une matrice de rotation en angles 
        d'Euler (pitch, yaw, roll).
        
        Formules:
        pitch = atan2(-R31, sqrt(R32^2 + R33^2))
        yaw   = atan2(R21, R11)
        roll  = atan2(R32, R33)
        """
        sy = np.sqrt(R[0, 0]**2 + R[1, 0]**2)
        
        if sy > 1e-6:  # Non singulier
            pitch = np.arctan2(-R[2, 0], sy)
            yaw = np.arctan2(R[1, 0], R[0, 0])
            roll = np.arctan2(R[2, 1], R[2, 2])
        else:  # Singulier (gimbal lock)
            pitch = np.arctan2(-R[2, 0], sy)
            yaw = np.arctan2(-R[1, 2], R[1, 1])
            roll = 0.0
        
        return np.array([pitch, yaw, roll])
\end{lstlisting}

\section{Implémentation de l'Analyse Comportementale}

\begin{lstlisting}[style=pythonstyle, caption={Module d'analyse comportementale.}, label={lst:behavioral}]
from collections import deque
import time

class BehavioralAnalyzer:
    """
    Analyse comportementale du conducteur.
    
    Calcule les indicateurs de fatigue et distraction
    sur une fenetre temporelle glissante.
    """
    
    def __init__(self, window_seconds=60, fps=30):
        """
        Args:
            window_seconds: Taille de la fenetre (secondes)
            fps: Images par seconde
        """
        self.window_size = window_seconds * fps
        self.fps = fps
        
        # Tampons circulaires
        self.eye_states = deque(maxlen=self.window_size)
        self.timestamps = deque(maxlen=self.window_size)
        
        # Suivi des clignements
        self.blink_events = deque(maxlen=100)
        self.current_blink_start = None
        self.prev_eye_state = 'open'
        
        # Indicateurs courants
        self.perclos = 0.0
        self.blink_frequency = 0.0
        self.avg_blink_duration = 0.0
        self.current_closure_duration = 0.0
    
    def update(self, eye_state, head_pose, timestamp=None):
        """
        Met a jour l'analyse avec une nouvelle observation.
        
        Args:
            eye_state: 'open' ou 'closed'
            head_pose: dict {'yaw', 'pitch', 'roll'}
            timestamp: Temps courant (secondes)
        """
        if timestamp is None:
            timestamp = time.time()
        
        self.eye_states.append(eye_state)
        self.timestamps.append(timestamp)
        
        # Mise a jour des clignements
        self._update_blinks(eye_state, timestamp)
        
        # Calcul des indicateurs
        self._compute_perclos()
        self._compute_blink_metrics()
        self._compute_closure_duration(eye_state, timestamp)
        
        # Calcul des scores
        fatigue_score = self._compute_fatigue_score()
        distraction_score = self._compute_distraction_score(
            head_pose
        )
        
        return {
            'perclos': self.perclos,
            'blink_frequency': self.blink_frequency,
            'avg_blink_duration': self.avg_blink_duration,
            'fatigue_score': fatigue_score,
            'distraction_score': distraction_score,
            'driver_state': self._classify_state(
                fatigue_score, distraction_score
            )
        }
    
    def _update_blinks(self, eye_state, timestamp):
        """Detecte et enregistre les clignements."""
        if self.prev_eye_state == 'open' and \
           eye_state == 'closed':
            # Debut d'un clignement
            self.current_blink_start = timestamp
        
        elif self.prev_eye_state == 'closed' and \
             eye_state == 'open':
            # Fin d'un clignement
            if self.current_blink_start is not None:
                duration = timestamp - self.current_blink_start
                if 0.05 < duration < 2.0:  # Filtre
                    self.blink_events.append({
                        'time': timestamp,
                        'duration': duration
                    })
                self.current_blink_start = None
        
        self.prev_eye_state = eye_state
    
    def _compute_perclos(self):
        """Calcule le PERCLOS sur la fenetre courante."""
        if len(self.eye_states) == 0:
            self.perclos = 0.0
            return
        
        closed_count = sum(
            1 for s in self.eye_states if s == 'closed'
        )
        self.perclos = (closed_count / len(self.eye_states)) \
                       * 100.0
    
    def _compute_fatigue_score(self):
        """
        Calcule le score de fatigue (equation 3.1).
        
        S_fatigue = 0.4*C_PERCLOS + 0.2*C_blink 
                  + 0.2*C_duration + 0.2*C_closure
        """
        # Composante PERCLOS
        c_perclos = min(self.perclos / 100.0, 1.0)
        
        # Composante frequence de clignement
        normal_freq = 17.5  # clig/min
        if self.blink_frequency > 0:
            c_blink = min(
                abs(self.blink_frequency - normal_freq) 
                / normal_freq, 1.0
            )
        else:
            c_blink = 0.5  # Absence de clignements
        
        # Composante duree des clignements
        c_duration = min(
            self.avg_blink_duration / 0.5, 1.0
        )
        
        # Composante fermeture courante
        c_closure = min(
            self.current_closure_duration / 2.0, 1.0
        )
        
        # Score combine
        score = (0.4 * c_perclos + 0.2 * c_blink + 
                 0.2 * c_duration + 0.2 * c_closure)
        
        return min(max(score, 0.0), 1.0)
    
    def _compute_distraction_score(self, head_pose):
        """
        Calcule le score de distraction (equation 3.2).
        
        S_distraction = 0.4*C_yaw + 0.2*C_pitch 
                      + 0.4*C_away
        """
        yaw = abs(head_pose.get('yaw', 0.0))
        pitch = abs(head_pose.get('pitch', 0.0))
        
        c_yaw = min(yaw / 45.0, 1.0)
        c_pitch = min(pitch / 35.0, 1.0)
        
        # Regard detourne
        is_away = yaw > 30.0 or pitch > 25.0
        c_away = 1.0 if is_away else 0.0
        
        score = 0.4 * c_yaw + 0.2 * c_pitch + 0.4 * c_away
        
        return min(max(score, 0.0), 1.0)
    
    def _classify_state(self, fatigue_score, 
                         distraction_score):
        """
        Classifie l'etat du conducteur.
        
        Algorithm 2 du rapport.
        """
        if fatigue_score > 0.7:
            return 'DROWSY'
        elif fatigue_score > 0.4:
            return 'FATIGUED'
        elif distraction_score > 0.4:
            return 'DISTRACTED'
        else:
            return 'ATTENTIVE'
\end{lstlisting}

\section{Implémentation de la Logique ECU}

\begin{lstlisting}[style=pythonstyle, caption={Module de logique décisionnelle ECU.}, label={lst:ecu}]
from enum import IntEnum

class AlertLevel(IntEnum):
    """Niveaux d'alerte du systeme ECU."""
    NONE = 0
    WARNING = 2
    CRITICAL = 3
    EMERGENCY = 4

class ECUDecisionModule:
    """
    Module de logique decisionnelle simulant un ECU ADAS.
    
    Implemente une machine a etats avec 4 niveaux d'alerte
    et des conditions de transition hysteretiques.
    """
    
    def __init__(self):
        self.current_level = AlertLevel.NONE
        self.alert_history = []
        self.state_duration = 0.0
        self.last_update_time = None
    
    def update(self, fatigue_score, distraction_score, 
               timestamp=None):
        """
        Met a jour l'etat de l'ECU.
        
        Args:
            fatigue_score: Score de fatigue [0, 1]
            distraction_score: Score de distraction [0, 1]
            timestamp: Temps courant
            
        Returns:
            dict: {'alert_level': int, 
                   'action': str, 
                   'message': str}
        """
        S_f = fatigue_score
        S_d = distraction_score
        S_max = max(S_f, S_d)
        
        prev_level = self.current_level
        
        # Logique de transition
        if self.current_level == AlertLevel.NONE:
            if S_f > 0.35 or S_d > 0.35:
                self.current_level = AlertLevel.WARNING
        
        elif self.current_level == AlertLevel.WARNING:
            if S_f > 0.6 or S_d > 0.6:
                self.current_level = AlertLevel.CRITICAL
            elif S_f < 0.2 and S_d < 0.2:
                self.current_level = AlertLevel.NONE
        
        elif self.current_level == AlertLevel.CRITICAL:
            if S_f > 0.75:
                self.current_level = AlertLevel.EMERGENCY
            elif S_f < 0.5 and S_d < 0.5:
                self.current_level = AlertLevel.WARNING
        
        elif self.current_level == AlertLevel.EMERGENCY:
            if S_f < 0.6:
                self.current_level = AlertLevel.CRITICAL
        
        # Generation de l'action
        action = self._get_action()
        
        # Enregistrement
        if self.current_level != prev_level:
            self.alert_history.append({
                'time': timestamp,
                'from': prev_level,
                'to': self.current_level,
                'fatigue': S_f,
                'distraction': S_d
            })
        
        return {
            'alert_level': int(self.current_level),
            'action': action,
            'message': self._get_message(),
            'changed': self.current_level != prev_level
        }
    
    def _get_action(self):
        """Retourne l'action associee au niveau courant."""
        actions = {
            AlertLevel.NONE: 'none',
            AlertLevel.WARNING: 'visual_alert',
            AlertLevel.CRITICAL: 'audio_alert',
            AlertLevel.EMERGENCY: 'emergency_brake'
        }
        return actions[self.current_level]
    
    def _get_message(self):
        """Retourne le message pour le tableau de bord."""
        messages = {
            AlertLevel.NONE: 
                'Systeme DMS: Fonctionnement normal',
            AlertLevel.WARNING: 
                'ATTENTION: Signes de fatigue detectes',
            AlertLevel.CRITICAL: 
                'ALERTE: Fatigue significative!',
            AlertLevel.EMERGENCY: 
                'URGENCE: Freinage d\'urgence active!'
        }
        return messages[self.current_level]
\end{lstlisting}

\section{Fonction de Perte et Métriques}

\subsection{Matrice de Confusion}

La performance du CNN est évaluée par la matrice de confusion :

\begin{equation}
\text{Matrice} = \begin{pmatrix} VP & FP \\ FN & VN \end{pmatrix}
\label{eq:confusion_matrix}
\end{equation}

Les métriques dérivées sont :

\begin{align}
\text{Précision} &= \frac{VP}{VP + FP} \label{eq:precision} \\
\text{Rappel (Sensibilité)} &= \frac{VP}{VP + FN} \label{eq:recall} \\
\text{Spécificité} &= \frac{VN}{VN + FP} \label{eq:specificity} \\
F_1 &= 2 \cdot \frac{\text{Précision} \cdot \text{Rappel}}{\text{Précision} + \text{Rappel}} = \frac{2 \cdot VP}{2 \cdot VP + FP + FN} \label{eq:f1}
\end{align}

\section{Synthèse}

Ce chapitre a présenté l'implémentation complète du système DMS, couvrant :
\begin{itemize}
    \item L'acquisition vidéo avec prétraitement CLAHE
    \item Deux architectures CNN (FaceDetectionCNN et EyeStateCNN) avec leurs processus d'entraînement
    \item La détection de visage/yeux par combinaison Haar + CNN
    \item L'estimation de la pose de la tête par PnP
    \item L'analyse comportementale avec calcul des scores de fatigue et distraction
    \item La logique décisionnelle ECU avec machine à états
\end{itemize}

L'ensemble du code est structuré en modules Python indépendants, testables et réutilisables.
